\documentclass[11pt]{article}
\usepackage{amsmath,amssymb}
\newcommand{\Gen}[1]{\langle #1\rangle}
\title{Two finite closure certificates in the original job 155 domain}
\author{}
\date{}
\begin{document}
\maketitle

The source census is Section 4 of Malcolm Hoong Wai Chen,
\emph{Abstract regular polytopes of $G_2(q)$ in even characteristic},
arXiv:2609.25004v1 (2026). The present positions refer to our finite
verification domain; they are not theorem numbers in that paper.
We borrow the accepted faithful natural matrix model of the original
$G_2(11)$, its original orbit catalogue, and its transporters. In particular,
the inequality $|G_2(11)|>32768$ is an explicit ambient premise, not an
identification made by the finite checker.

\paragraph{Proposition.}
The original representatives at positions $328$ and $682$ of job $[155,4]$
do not generate the intended ambient group. Their catalogue indices are
respectively $3152$ and $8052$. In each case the four actual generators
generate a group of order $12096$.

\paragraph{Proof.}
Write the four involutions as $r_0,r_1,r_2,r_3$. The supplied certificate
is a rooted tree with $12096$ vertices. Its root represents the identity;
every other vertex has a strictly earlier parent and a label from
$\{0,1,2,3\}$. Recursively, if vertex $v$ has parent $p$ and label $j$,
its matrix is $T_v=T_p r_j$. Thus every $T_v$ lies in
$H=\Gen{r_0,r_1,r_2,r_3}$.

The separate verifier reconstructs these matrices from the tree and the
four original generators. It checks that all $12096$ matrices are distinct.
It then checks each of the $4\cdot12096=48384$ products $T_vr_j$ against
this finite set, comparing all matrix entries. Put $E=\{T_v\}$. The root
gives $1\in E$, and the checked products give $Er_j\subseteq E$ for each
$j$. Every word in the four generators consequently belongs to $E$.
Since the generators are involutions, these words constitute $H$.
We have both $H\subseteq E$ and $E\subseteq H$, so
\[
  |H|=|E|=12096<32768<|G_2(11)|.
\]
The borrowed faithful representation and accepted original transporters
carry this nongeneration statement to the specified original orbit entries.
\hfill$\square$

\paragraph{Verification and attribution.}
The proof uses the elementary closure criterion for a finitely generated
group. Its tree description stores actual words; it does not infer group
membership from membership in a matrix algebra. The new work is the two
supplied finite certificates, their separate complete checks, and this
application to the recorded original representatives. The model,
identification of the original domain, and ambient order premise are
borrowed from the previously accepted inputs.

The Source packet is \texttt{f004\_finite\_closure}; its independent static
review and twenty-two fresh operational controls precede the Original
calls. The case reviews \texttt{f004\_job155\_case328\_closure} and
\texttt{f004\_job155\_case682\_closure} bind the exact Source, full input
frontier, original case, genuine producer and checker calls, and resource
receipts. Their role is to review those executions and their logical
interpretation, not to claim an additional numerical run.

These are two negative census entries. They neither close the job nor
identify the smaller group, decide an open conjecture, establish the
intersection property, or supply a Lean formalization.
\end{document}
